\chapter{Standard Strategies}
\begin{quote}
  Uh, perpendiculars should be okay right?  Dot product?
  \\ -- Evan Chen, on EFFT
\end{quote}

This chapter will cover some typical techniques which arise in olympiad problems.

First, we begin the following definition \cite{infinity}:
\begin{StoreCode}{tech}
\begin{definition}
  The \textit{displacement vector} of two (normalized) points $P = (p_1, p_2, p_3)$ and $Q = (q_1, q_2, q_3)$ is denoted by $\overrightarrow{PQ}$ and is equal to $(p_1-q_1, p_2-q_2, p_3-q_3)$.
\end{definition}
\end{StoreCode}

Note that the sum of the coordinates of a displacement vector is $0$.

In the proofs of this chapter, we will invoke some properties which will be familiar to those who have done vector bashing; namely, when $\vec O = \vec 0$
\begin{align*}
  \vec A \cdot \vec A &= R^2 \\
  \vec A \cdot \vec B &= R^2-c^2/2
\end{align*}
where $R$ is the circumradius.

If you don't care about the proofs, then this doesn't matter; otherwise, see appendix \ref{appendix:vectors} for a treatise on some elementary vector bashing.

This chapter will frequently involve translating $O$ to the null vector $\vec 0$, so that we may invoke useful properties about dot products mentioned above.  This is valid since the point $(x,y,z)$ satisfies $x+y+z=1$, so translation is actually an invariant.  As a consequence, however, take heed of the following.
\begin{note}
  It is important that $x+y+z=1$ when doing calculations!
\end{note}





\StoreText{\subsubsection*{A Note on Scaling Displacement Vectors} In EFFT, one can write a displacement vector $(x,y,z)$ as $(kx:ky:kz)$, and the theorem will still be true.  This is also true for Strong EFFT, but NOT for the distance formula.}{tech}

\section{EFFT: Perpendicular Lines}
In fact there is a fast characterization for when two lines $MN$ and $PQ$ are perpendicular.  This lemma also appears in \cite{book}, with a proof via the Pythagorean Theorem.
\begin{StoreCode}{tech}
\begin{theorem}[Evan's Favorite Forgotten Trick]
  \label{tech:efft}
  Consider displacement vectors $\overrightarrow{MN} = (x_1, y_1, z_1)$ and $\overrightarrow{PQ} = (x_2, y_2, z_2)$.  Then $MN \perp PQ$ if and only if
  \[  0 = a^2(z_1y_2 + y_1z_2) + b^2(x_1z_2 + z_1x_2) + c^2(y_1x_2 + x_1y_2)\]
\end{theorem}
\end{StoreCode}
\begin{proof}
  Translate $\vec O$ to $\vec 0$.  It is necessary and sufficient that $(x_1\vec A + y_1\vec B + z_1 \vec C) \cdot (x_2 \vec A + y_2 \vec B + z_2 \vec C) = 0$.  Expanding, this is just
  \[ \cycsum \left(x_1x_2 \vec A \cdot \vec A \right) + \cycsum \left( \left(x_1y_2 + x_2y_1\right) \vec A \cdot \vec B \right) = 0 \]

  Utilizing the fact that $\vec O = \vec 0$, we make the standard substitutions and rearrange:

  \begin{align*}
    \iff \cycsum (x_1x_2 R^2) + \cycsum \left(x_1y_2 + x_2y_1\right)\left(R^2 - \frac{c^2}{2}\right) &= 0 \\
    \iff R^2 \left( \cycsum (x_1x_2) + \cycsum(x_1y_2 + x_2y_1) \right) &= \half \cycsum \left( (x_1y_2 + x_2y_1)(c^2) \right) \\
    \iff R^2 (x_1 + y_1 + z_1)(x_2 + y_2 + z_2) &= \half \cycsum \left( (x_1y_2 + x_2y_1)(c^2) \right) \\
    \iff R^2 \cdot 0 \cdot 0 &= \half \cycsum \left( (x_1y_2 + x_2y_1)(c^2) \right) \\
    \iff 0 &= \cycsum \left( (x_1y_2 + x_2y_1)(c^2) \right)
  \end{align*}
  since $x_1+y_1+z_1 = x_2+y_2+z_2 = 0$ (remember that $(x_1, y_1, z_1)$ is a displacement vector!).
\end{proof}
We will see later, in Section \ref{sec:notnormal}, that this holds for \textit{scaled} displacement vectors.

This will be abbreviated EFFT\footnote{Hi Simon.} throughout this document.  As we shall see, EFFT increases the scope of solvable problems by providing a usable method for constructing perpendiculars, and in particular, perpendicular bisectors.

EFFT is particularly useful when one of the displacement vectors is the side of the triangle.  For instance, we may deduce
\begin{StoreCode}{tech}
\begin{corollary}
  Consider a displacement vector $\overrightarrow{PQ} = (x_1,y_1,z_1)$.  Then $PQ \perp BC$ if and only if
  \[ 0 = a^2(z_1-y_1) + x_1(c^2-b^2) \]
\end{corollary}
\begin{corollary}
  \label{tech:perpbis}
  The perpendicular bisector of $BC$ has equation
  \[ 0 = a^2(z-y) + x(c^2-b^2) \]
\end{corollary}
\end{StoreCode}
\begin{exercise}
  Prove the above corollaries, and deduce the Pythagorean Theorem.
\end{exercise}

EFFT is used in many problems, e.g. problems \ref{exam:2008usamo2} and \ref{exam:2012wootpo47}.  Its generalization Strong EFFT (cf. Theorem \ref{tech:strongefft}) is also useful and is invoked in problem \ref{exam:2005islg5}.






\section{Distance Formula}
The distance formula \cite{infinity} can be proved by precisely the same means as EFFT:
\begin{StoreCode}{tech}
\begin{theorem}[Distance Formula]
  \label{tech:distance}
  Consider a displacement vector $\overrightarrow{PQ} = (x,y,z)$.  Then \[ |PQ|^2 = -a^2yz - b^2zx - c^2xy \]
\end{theorem}
\end{StoreCode}
\begin{exercise}
  Prove the distance formula by using \[ PQ^2 = (x\vec A + y\vec B + z \vec C) \cdot (x \vec A + y \vec B + z \vec C) \] and the fact that $x+y+z=0$.
\end{exercise}









\section{Circles}

\subsection{Equation of the Circle}
Not surprisingly, the equation of a circle \cite{infinity} is much more annoying than anything that has showed up so far.  Nonetheless, it is certainly usable, particularly in certain situations described below.

\begin{StoreCode}{tech}
\begin{theorem}
  \label{tech:circle}
  The general equation of a circle is \[ -a^2yz - b^2zx -c^2xy + (ux+vy+wz)(x+y+z) = 0 \] for reals $u,v,w$.
\end{theorem}
\end{StoreCode}
\begin{proof}
  Assume the circle has center $(j,k,l)$ and radius $r$.  Then this is just
  \[ -a^2(y-k)(z-l) - b^2(z-l)(x-j) - c^2(x-j)(y-k) = r^2 \]
  Expand everything, and collect terms to get
  \[ -a^2yz - b^2zx - c^2xy + C_1 x + C_2 y + C_3 z = C \]
  for some hideous constants $C_i$ and $C$.  Since $x+y+z = 1$, we can rewrite this as
  \[ -a^2yz - b^2zx - c^2xy + ux + vy + wz = 0 \iff -a^2xy - b^2yz -c^2zx + (ux+vy+wz)(x+y+z) = 0\]
  where $u=C_1-C$, et cetera.
\end{proof}

Fortunately, most circles that we are concerned with have much cleaner forms.  For example, as in \cite{infinity},
\begin{StoreCode}{tech}
\begin{corollary}
  The circumcircle has equation \[ a^2yz + b^2zx + c^2xy = 0 \]
\end{corollary}
\end{StoreCode}
\begin{exercise}
  Prove the above corollary.
\end{exercise}

Furthermore, circles passing through vertices of the triangle make for more elegant circle formulas, since many terms in the circle equation vanish.

\begin{note}
  Notice that the equation of the circle is homogeneous, so it is possible to scale $(x,y,z)$ arbitrarily and still have a true equation.  See section \ref{sec:notnormal}.
\end{note}
